\section{Code Review 体系}

\subsection{Review 的时机与原则}

\begin{importantbox}{{Review Early, Review Often}}
\begin{itemize}
    \item 不是在合并主分支前才审查，而是\textbf{每个细粒度子任务完成后}都必须强制触发审查。
    \item 防止设计错误"污染"下一个任务。
    \item Code Review 的输入：用户原始需求/计划 + Git Diff。
\end{itemize}
\end{importantbox}

\subsection{请求 Code Review 的 Checklist}

\begin{enumerate}
    \item \textbf{代码质量}：关注点分离、解耦、错误处理、重复代码、边界 Case 覆盖。
    \item \textbf{架构}：设计角色是否合理、可扩展性、性能影响、安全隐患。
    \item \textbf{测试}：是否测了真实逻辑（而非只是 Mock）、边界覆盖率、集成测试。
    \item \textbf{需求}：是否满足计划的所有要求、无功能蔓延、达到生产级完成度。
\end{enumerate}

\begin{warningbox}{行为红线}
\begin{itemize}
    \item 严禁不看代码就说"看起来不错"。
    \item 严禁将次要问题标为致命错误。
    \item 严禁含糊其词——必须给出最终结论。
\end{itemize}
\end{warningbox}

\subsection{接收 Code Review 的准则}

Code Review 的本质是\textbf{基于客观代码的技术评估}，而非社交表现：

\begin{itemize}
    \item \textbf{追求技术严谨性}：决策基于事实、测试结果、系统架构现状，而非直觉或权威。
    \item \textbf{拒绝表演性同意}：不在未验证技术可行性的前提下盲目同意 Reviewer 建议。
    \item \textbf{YAGNI 原则}：若 Reviewer 建议添加当前系统根本没有调用的功能，应拒绝。
    \item \textbf{验证先于执行}：收到反馈后第一步是去代码库中验证，而非立即修改。
    \item \textbf{有理有据的技术反驳}（Push Back）：发现 Reviewer 建议有误时，使用技术推理反驳——反驳不是为了争吵，而是为了保护代码库。
\end{itemize}

\begin{knowledgebox}{消除模糊性}
如果 Reviewer 提了 6 点意见，其中 2 点描述不清，\textbf{严禁先做懂的 4 点、把不懂的 2 点留到最后}。必须立刻停下来，先询问清楚所有细节——因为软件模块之间高度耦合，部分理解会导致错误的实现方向，返工成本极高。
\end{knowledgebox}

\begin{importantbox}{行动胜于雄辩}
当 Reviewer 意见正确时，严禁使用感激或赞美的废话。正确做法是\textbf{直接修复并简短陈述事实}。高效工程文化中，代码本身就是最好的回应。
\end{importantbox}

\subsection{SOLID 与关注点分离}

Code Review 中涉及的架构设计模式：

\begin{enumerate}
    \item \textbf{单一职责原则}（SRP）：避免一个类/文件承担所有事情，必须解耦。
    \item \textbf{开闭原则}（OCP）：对扩展开放、对修改关闭。
    \item \textbf{里氏替换原则}（LSP）：子类可替换父类。
    \item \textbf{接口隔离原则}（ISP）：接口应细粒度。
    \item \textbf{依赖倒置原则}（DIP）：依赖抽象而非具体实现。
\end{enumerate}

\textbf{关注点分离}是宏观架构哲学：将软件划分为具有不同职责的独立模块，每个模块只解决一个特定的关注点。

\subsection{本章小结}

Code Review 体系包含"请求审查"和"接收审查"两个 Skill，通过 Checklist 和行为红线确保审查的技术严谨性，同时用 SOLID 原则和关注点分离指导架构评估。


